This subsection asks who pays the bias measured in \cref{sec:method-xs}. It uses every public
print in the trade sample of \cref{sec:data-flow}. Kalshi's public trade feed has no account
identities, but it does name the side that crossed the spread. That is enough to separate the two
parties to each trade by \emph{liquidity role}.

\paragraph{Takers, makers and positions.} A print $i$ records a number of contracts $n_i$, a YES
price $p_i$ and the taker's side.
\begin{itemize}
  \item A taker who bought YES holds YES at $\q_i=p_i$. A taker who bought NO holds NO at
  $\q_i=1-p_i$.
  \item The resting order on the other side (the \emph{maker}) holds the opposite side at
  $1-\q_i$.
  \item Gross of fees, the two positions are mirror images. The taker earns $\won_i-\q_i$ per
  contract and the maker earns $-(\won_i-\q_i)$.
  \item Fees break the symmetry. The taker pays $\fee^T_i=0.07\,\kappa\,\q_i(1-\q_i)$, where
  $\kappa$ is the series multiplier. The maker pays $\fee^M_i=0.0175\,\kappa\,\q_i(1-\q_i)$ on
  series that charge makers, and nothing otherwise.
\end{itemize}
Each print therefore contributes two \emph{positions}, one per role. We sort positions by the
price of the side held, $\q$ for the taker and $1-\q$ for the maker, into the buckets of
\cref{sec:method}.

\paragraph{Weights and returns.} Let $\omega_i$ be the design weight of the print's market
(the inverse probability that the sampling design of \cref{sec:data-flow} drew it). For any cell
of positions we report two averages.
\begin{align}
  \bar\Delta &= \frac{\sum_i \omega_i n_i(\won_i-\q_i-\fee_i)}{\sum_i\omega_i n_i}
  &&\text{(contract-weighted, \$ per contract),}\label{eq:flow-c}\\
  \bar\ret &= \frac{\sum_i \omega_i n_i(\won_i-\q_i-\fee_i)}{\sum_i\omega_i n_i(\q_i+\fee_i)}
  &&\text{(dollar-weighted, per dollar).}\label{eq:flow-d}
\end{align}
Here $\q_i$, $\won_i$ and $\fee_i$ refer to the position's own side. \Cref{eq:flow-d} is the
dollar-weighted mean of the per-dollar return~\eqref{eq:ret}, that is, aggregate profit over
aggregate stake. Unless a table says otherwise, returns are net of the exact fee. Gross returns
set $\fee_i=0$. As a retail-perspective variant we also report a print-weighted mean, in which
each fill counts once regardless of its size.

\paragraph{Taker orders and trade size.} Kalshi prints one fill for each resting order that a
taker order executes against. A taker order that sweeps several resting orders therefore
appears as several prints. These prints share the market, the microsecond timestamp and the
taker side. We group such prints into one \emph{taker order} $o$, with dollar notional
$N_o=\sum_{i\in o} n_i\q_i$. Trade size is $N_o$, bucketed as $\le\$10$, $\$10$--$100$,
$\$100$--$1\text{k}$ and $>\$1\text{k}$. Every print inherits the size of its order. Small trade
size is the standard proxy for retail flow when identities are hidden
\citep{leeradhakrishna2000,barberodeanzhu2009}. It is only a proxy: a professional can split
orders, and a retail trader can send a large one. Makers' order sizes are not observed. A
print shows only the part of the resting order that was filled.

\paragraph{Timing and attention.} Each print is classified by the following variables. Each is
known at the time of the print, except the marquee label.
\begin{itemize}
  \item \emph{Clock.} Hour of day and weekday or weekend, in US Eastern time.
  \item \emph{In play.} A sports print is in play if it is at or after the game's scheduled
  start. The start is the time encoded in the ticker where there is one. Otherwise it is the
  occurrence time minus three hours, Kalshi's convention for game markets. A start is used only
  if it falls between 30 minutes and 12 hours before the close.
  \item \emph{Primetime.} The game is scheduled to start between 19:00 and 23:59 ET.
  \item \emph{Time to expiry.} Hours from the print to the market's scheduled expiration.
  \item \emph{New listing.} The print falls in the first 24 hours after the market opened.
  \item \emph{Volume surge.} The market's volume in the hour before the print, $V_{t-1h,t}$, is
  compared with its own average hourly volume before that hour, $\bar V_{t-1h}$. Both use past
  prints only, and the average needs at least six hours of history. A surge is
  $V_{t-1h,t}\ge3\,\bar V_{t-1h}$, which captures attention shocks in the spirit of
  \citet{barberodean2008}.
  \item \emph{Marquee.} The event's total volume is in the top decile of its stratum. This label
  is descriptive and uses ex post volume.
\end{itemize}

\paragraph{Mid-quote proxy, price impact and markouts.} The public trade history contains no
quotes. We therefore build a mid-quote from the prints themselves.
\begin{itemize}
  \item A YES-taker print executes at the YES ask, and a NO-taker print executes at the YES bid.
  \item Let $a_t$ be the YES price of the market's last YES-taker print at or before $t$, and
  $b_t$ that of its last NO-taker print.
  \item The proxy is $m_t=(a_t+b_t)/2$. It is defined only when both prints are at most 60
  minutes old and $a_t\ge b_t$.
  \item Impact, markout and order-flow analyses also require a tight proxy,
  $a_t-b_t\le5\cents$, both before the order and at each later horizon. This keeps them to
  markets where the proxy tracks the book.
\end{itemize}
For a taker order at time $t$ with fill price $p$ and direction $d=+1$ (bought YES) or $d=-1$
(bought NO), define, in cents per contract,
\begin{equation}\label{eq:flow-markout}
  s=d\,(p-m_{t^-}),\qquad I_h=d\,(m_{t+h}-m_{t^-}),\qquad M_h=d\,(m_{t+h}-p)=I_h-s,
\end{equation}
where $m_{t^-}$ is the proxy just before the order. These are the effective half-spread $s$,
the price impact $I_h$ and the taker's markout $M_h$ at horizons $h$ of one minute and one hour.
A mark at $t+h$ is used only if $t+h$ is before the market's pre-resolution cut, so no markout
reads a price set after the outcome was known. The markout to settlement, $M_\infty=d\,(Y-p)$
with $Y$ the YES outcome, is the taker's gross profit per contract. It splits into the
mispricing of the mid and the cost of immediacy, $M_\infty=d\,(Y-m_{t^-})-s$. Persistent impact
signals information, and impact that reverses signals inventory or noise
\citep{glostenmilgrom1985,hasbrouck1991}.

We validated the proxy against order books recorded from Kalshi's WebSocket feed in two
three-hour windows.
% source: data/research/flb_flow/l2check.json
\begin{itemize}
  \item \emph{Saturday evening ET (8~November, 00--03~UTC).} The window covers 236 markets and
  193,907 taker orders. 98\% of its prints are in game-winner markets, mainly college football
  and NBA games during the evening slate. The proxy exists for 96\% of orders. It is within one
  cent of the book mid for 96.9\% of them and within two cents for 98.9\%. Its mean absolute
  error is 0.16\cents.
  \item \emph{Tuesday morning ET (10~November, 14--17~UTC).} The window covers 174 markets and
  11,717 orders. It is mainly pre-game NBA and NHL markets, with daily temperature, mentions and
  political markets. The proxy exists for 86\% of orders, is within one cent for 94.4\% and
  within two cents for 97.0\%. Its mean absolute error is 0.27\cents.
\end{itemize}
In both windows the proxy's correlation with the book mid exceeds 0.998, and 99.1\% of taker
orders execute at or beyond the book mid on their own side. This confirms the meaning of the
taker-side field. The recorded books are among the more liquid on the venue, with a median
spread of one cent. In thin markets the proxy is staler and more often missing.

\paragraph{Order-flow regressions.} For the information content of flow we aggregate taker
orders to market-hours $h$, before each market's pre-resolution cut. Order sizes fall into three
classes: small ($\le\$100$), medium ($\$100$--$1$k) and large ($>\$1$k). Class $c$'s signed
imbalance is
\begin{equation}\label{eq:flow-oib}
  \mathrm{OIB}_{c,h}=\frac{\sum_{o\in(c,h)} d_o n_o}{\sum_{o\in h} n_o}\in[-1,1].
\end{equation}
It is normalised by the hour's total volume, so the three classes sum to the hour's net taker
imbalance. We then estimate
\begin{equation}\label{eq:flow-reg}
  \Delta^{(k)}_{h}=\alpha+\sum_{c}\beta^{(k)}_c\,\mathrm{OIB}_{c,h}
  +\gamma^{(k)}\,(m^{\mathrm{end}}_h-m^{\mathrm{start}}_h)+\varepsilon_h .
\end{equation}
The left-hand side $\Delta^{(k)}_h$ is measured in cents and takes three forms: the same-hour mid
change $m^{\mathrm{end}}_h-m^{\mathrm{start}}_h$ (with $\gamma=0$), the next-hour change, and the
drift to settlement, $Y-m^{\mathrm{end}}_h$. A class whose flow moves the price and is
confirmed at settlement ($\beta^{(\infty)}_c>0$) carries information in the sense of
\citet{kyle1985}. A class whose same-hour impact reverses ($\beta^{(\mathrm{next})}_c<0$,
$\beta^{(\infty)}_c\le0$) behaves as noise. The settlement drift needs no mid at the outcome,
because $Y$ is observed.

\paragraph{Inference.} Standard errors are clustered by event \citep{cameronmiller2015}. For
weighted means we compute them from per-event sums of $\omega_in_i(x_i-\bar x)$, which gives
exactly the cluster-robust variance. A difference between two cells, such as takers against
makers or YES-takers against NO-takers, uses the per-event influence functions of both means.
This accounts for the two cells sharing events. Regression~\eqref{eq:flow-reg} uses CR1
cluster-robust errors. Every table reports $t$-statistics and the number of event clusters.
